\documentclass{beamer}
\usepackage[utf8]{inputenc}
\usepackage[spanish]{babel}
\usepackage{listings}
\usepackage{xcolor}
\usepackage{tikz}
\usepackage{tcolorbox}
\usepackage{fontawesome5}

% Configuración de colores - Gruvbox Light con marrón principal
\definecolor{bg0}{HTML}{fbf1c7}
\definecolor{bg1}{HTML}{ebdbb2}
\definecolor{bg2}{HTML}{d5c4a1}
\definecolor{bg3}{HTML}{bdae93}
\definecolor{fg0}{HTML}{282828}
\definecolor{fg1}{HTML}{3c3836}
\definecolor{fg2}{HTML}{504945}
\definecolor{aqua}{HTML}{689d6a}
\definecolor{aquabright}{HTML}{427b58}
\definecolor{blue}{HTML}{458588}
\definecolor{red}{HTML}{cc241d}
\definecolor{orange}{HTML}{d65d0e}
\definecolor{yellow}{HTML}{d79921}
\definecolor{green}{HTML}{98971a}
\definecolor{purple}{HTML}{b16286}
\definecolor{brown}{HTML}{AB886D}
\definecolor{brownbright}{HTML}{8B6849}

% Colores para código
\definecolor{codegreen}{HTML}{98971a}
\definecolor{codegray}{HTML}{7c6f64}
\definecolor{codepurple}{HTML}{b16286}
\definecolor{backcolour}{HTML}{f9f5d7}

% Configuración de código
\lstdefinestyle{pythonstyle}{
    backgroundcolor=\color{backcolour},   
    commentstyle=\color{codegreen},
    keywordstyle=\color{brownbright}\bfseries,
    numberstyle=\tiny\color{codegray},
    stringstyle=\color{codepurple},
    basicstyle=\ttfamily\footnotesize\color{fg0},
    breakatwhitespace=false,         
    breaklines=true,                 
    captionpos=b,                    
    keepspaces=true,                 
    numbers=left,                    
    numbersep=3pt,                  
    showspaces=false,                
    showstringspaces=false,
    showtabs=false,                  
    tabsize=2,
    language=Python,
    frame=single,
    rulecolor=\color{brown},
    xleftmargin=8pt,
    xrightmargin=8pt,
    inputencoding=utf8,
    extendedchars=true,
    literate=
        {á}{{\'a}}1 {é}{{\'e}}1 {í}{{\'i}}1 {ó}{{\'o}}1 {ú}{{\'u}}1
        {Á}{{\'A}}1 {É}{{\'E}}1 {Í}{{\'I}}1 {Ó}{{\'O}}1 {Ú}{{\'U}}1
        {ñ}{{\~n}}1 {Ñ}{{\~N}}1
        {¿}{{?`}}1 {¡}{{!`}}1
}

\lstset{style=pythonstyle}

% Tema personalizado con marrón principal
\usetheme{Madrid}
\setbeamercolor{structure}{fg=brown}
\setbeamercolor{palette primary}{bg=brown,fg=bg0}
\setbeamercolor{palette secondary}{bg=brownbright,fg=bg0}
\setbeamercolor{palette tertiary}{bg=fg1,fg=bg0}
\setbeamercolor{palette quaternary}{bg=brown,fg=bg0}
\setbeamercolor{block title}{bg=brown,fg=bg0}
\setbeamercolor{block body}{bg=bg1,fg=fg0}
\setbeamercolor{frametitle}{bg=brown,fg=bg0}
\setbeamercolor{title}{fg=fg0}
\setbeamercolor{subtitle}{fg=fg1}
\setbeamercolor{normal text}{fg=fg0,bg=bg0}
\setbeamercolor{item}{fg=brown}
\setbeamercolor{section in toc}{fg=brown}

% Cajas personalizadas con marrón
\newtcolorbox{conceptbox}[1]{
    colback=bg1,
    colframe=brown,
    fonttitle=\bfseries\small,
    title=#1,
    coltitle=fg0,
    left=3pt,
    right=3pt,
    top=3pt,
    bottom=3pt,
    boxsep=2pt
}

\newtcolorbox{notebox}{
    colback=yellow!15,
    colframe=orange,
    leftrule=2mm,
    rightrule=0mm,
    toprule=0mm,
    bottomrule=0mm,
    arc=0mm,
    left=3pt,
    right=3pt,
    top=2pt,
    bottom=2pt
}

\title{\textbf{Introducción a la Programación con Python}}
\subtitle{Semana 6: Entrada/Salida de Archivos}
\author{$\nabla$Nablabs}
\date{}
\begin{document}

\begin{frame}
\titlepage
\end{frame}

\begin{frame}
\frametitle{Contenido}
\tableofcontents
\end{frame}

\section{¿Por qué trabajar con archivos?}

\begin{frame}{Archivos: Persistencia de datos \faFileAlt}
\begin{conceptbox}{¿Por qué son importantes?}
\small
\begin{itemize}
    \item[\faSave] Guardar datos permanentemente
    \item[\faDatabase] Almacenar configuraciones y logs
    \item[\faExchangeAlt] Intercambiar información entre programas
    \item[\faChartLine] Procesar grandes volúmenes de datos
    \item[\faHistory] Mantener historial de operaciones
\end{itemize}
\end{conceptbox}

\vspace{0.2cm}
\begin{center}
\large\textcolor{brown}{\faLightbulb\ Los archivos permiten que tus programas "recuerden"}
\end{center}
\end{frame}

\section{Función open()}

\begin{frame}[fragile]{open(): Abrir archivos \faFolderOpen}
\begin{conceptbox}{¿Qué hace open()?}
Abre un archivo y retorna un objeto de archivo para manipularlo
\end{conceptbox}

\textbf{Sintaxis básica:}
\begin{lstlisting}
archivo = open("nombre.txt", "modo")
# ... operaciones con el archivo ...
archivo.close()  # ¡Siempre cerrar!
\end{lstlisting}

\vspace{0.2cm}
\textbf{Componentes:}
\small
\begin{itemize}
    \item[\faFile] \textbf{nombre.txt}: Ruta del archivo
    \item[\faKey] \textbf{modo}: Cómo abrir el archivo (lectura, escritura, etc.)
    \item[\faTimes] \textbf{close()}: Cierra el archivo liberando recursos
\end{itemize}

\vspace{-0.1cm}
\begin{notebox}
\small\faExclamationTriangle\ Olvidar cerrar un archivo puede causar problemas
\end{notebox}
\end{frame}

\section{Modos de apertura}

\begin{frame}[fragile]{Modos de apertura \faKey}
\begin{conceptbox}{Modos principales}
\small
\begin{itemize}
    \item[\faBookOpen] \textbf{'r'} (read): Solo lectura. Error si no existe.
    \item[\faPencilAlt] \textbf{'w'} (write): Escritura. Crea o sobreescribe.
    \item[\faPlus] \textbf{'a'} (append): Añade al final. Crea si no existe.
    \item[\faEdit] \textbf{'r+'}: Lectura y escritura. Error si no existe.
\end{itemize}
\end{conceptbox}

\vspace{0.2cm}
\textbf{Modos adicionales:}
\small
\begin{itemize}
    \item[\faFileCode] \textbf{'b'}: Modo binario (ej: 'rb', 'wb')
    \item[\faFileAlt] \textbf{'t'}: Modo texto (por defecto)
\end{itemize}

\vspace{-0.1cm}
\begin{notebox}
\small\faInfoCircle\ Por defecto, si no se especifica, usa 'r' (solo lectura)
\end{notebox}
\end{frame}

\begin{frame}[fragile]{Comparación de modos \faTable}
\begin{center}
\small
\begin{tabular}{|c|c|c|c|}
\hline
\textbf{Modo} & \textbf{Leer} & \textbf{Escribir} & \textbf{Si no existe} \\
\hline
\texttt{'r'} & \faCheck & \faTimes & Error \\
\hline
\texttt{'w'} & \faTimes & \faCheck & Crea archivo \\
\hline
\texttt{'a'} & \faTimes & \faCheck & Crea archivo \\
\hline
\texttt{'r+'} & \faCheck & \faCheck & Error \\
\hline
\end{tabular}
\end{center}

\vspace{0.3cm}
\begin{notebox}
\small\faExclamationTriangle\ \textbf{'w'} borra todo el contenido previo del archivo
\end{notebox}
\end{frame}

\section{Lectura de archivos}

\begin{frame}[fragile]{Leer archivo completo: read() \faBookReader}
\textbf{Archivo:} \texttt{saludo.txt}
\begin{tcolorbox}[colback=bg1,colframe=brown,left=3pt,right=3pt,top=2pt,bottom=2pt]
\small\texttt{Hola mundo\\
Bienvenido a Python\\
¡Archivos son útiles!}
\end{tcolorbox}

\vspace{0.2cm}
\begin{lstlisting}
# Leer todo el contenido
archivo = open("saludo.txt", "r")
contenido = archivo.read()
print(contenido)
archivo.close()
\end{lstlisting}

\vspace{-0.2cm}
\begin{tcolorbox}[colback=green!10,colframe=green,title={\small\faArrowCircleRight \ Salida},coltitle=fg0,left=3pt,right=3pt,top=2pt,bottom=2pt,boxsep=1pt]
\small\texttt{Hola mundo\\
Bienvenido a Python\\
¡Archivos son útiles!}
\end{tcolorbox}
\end{frame}

\begin{frame}[fragile]{Leer línea por línea: readline() \faStream}
\begin{lstlisting}
# Leer línea a línea
archivo = open("saludo.txt", "r")

linea1 = archivo.readline()
print(f"Línea 1: {linea1.strip()}")

linea2 = archivo.readline()
print(f"Línea 2: {linea2.strip()}")

archivo.close()
\end{lstlisting}

\vspace{-0.2cm}
\begin{tcolorbox}[colback=green!10,colframe=green,title={\small\faArrowCircleRight \ Salida},coltitle=fg0,left=3pt,right=3pt,top=2pt,bottom=2pt,boxsep=1pt]
\small\texttt{Línea 1: Hola mundo\\
Línea 2: Bienvenido a Python}
\end{tcolorbox}

\vspace{-0.1cm}
\begin{notebox}
\small\faInfoCircle\ \texttt{strip()} elimina espacios y saltos de línea
\end{notebox}
\end{frame}

\begin{frame}[fragile]{Leer todas las líneas: readlines() \faList}
\begin{lstlisting}
# Leer todas las líneas como lista
archivo = open("saludo.txt", "r")
lineas = archivo.readlines()
archivo.close()

print(f"Total de líneas: {len(lineas)}")
for i, linea in enumerate(lineas, 1):
    print(f"{i}. {linea.strip()}")
\end{lstlisting}

\vspace{-0.2cm}
\begin{tcolorbox}[colback=green!10,colframe=green,title={\small\faArrowCircleRight \ Salida},coltitle=fg0,left=3pt,right=3pt,top=2pt,bottom=2pt,boxsep=1pt]
\small\texttt{Total de líneas: 3\\
1. Hola mundo\\
2. Bienvenido a Python\\
3. ¡Archivos son útiles!}
\end{tcolorbox}

\vspace{-0.1cm}
\begin{notebox}
\small\faCheckCircle\ \texttt{readlines()} retorna una lista de strings
\end{notebox}
\end{frame}

\begin{frame}[fragile]{Iterar directamente sobre el archivo \faRedo}
\begin{lstlisting}
# Forma más eficiente y pythónica
archivo = open("saludo.txt", "r")

for linea in archivo:
    print(f">>> {linea.strip()}")

archivo.close()
\end{lstlisting}

\vspace{-0.2cm}
\begin{tcolorbox}[colback=green!10,colframe=green,title={\small\faArrowCircleRight \ Salida},coltitle=fg0,left=3pt,right=3pt,top=2pt,bottom=2pt,boxsep=1pt]
\small\texttt{>>> Hola mundo\\
>>> Bienvenido a Python\\
>>> ¡Archivos son útiles!}
\end{tcolorbox}

\vspace{-0.1cm}
\begin{notebox}
\small\faLightbulb\ Esta es la forma preferida: eficiente en memoria
\end{notebox}
\end{frame}

\section{Escritura en archivos}

\begin{frame}[fragile]{Escribir texto: write() \faPencilAlt}
\begin{lstlisting}
# Crear y escribir en un archivo
archivo = open("nuevo.txt", "w")

archivo.write("Primera línea\n")
archivo.write("Segunda línea\n")
archivo.write("Tercera línea\n")

archivo.close()

print("Archivo creado exitosamente")
\end{lstlisting}

\vspace{-0.2cm}
\begin{tcolorbox}[colback=green!10,colframe=green,title={\small\faArrowCircleRight \ Salida},coltitle=fg0,left=3pt,right=3pt,top=2pt,bottom=2pt,boxsep=1pt]
\small\texttt{Archivo creado exitosamente}
\end{tcolorbox}

\vspace{-0.1cm}
\begin{notebox}
\small\faExclamationTriangle\ Debes añadir \texttt{\textbackslash n} manualmente para saltos de línea
\end{notebox}
\end{frame}

\begin{frame}[fragile]{Escribir múltiples líneas: writelines() \faListAlt}
\begin{lstlisting}
# Escribir lista de líneas
frutas = ["Manzana\n", "Banana\n", "Cereza\n"]

archivo = open("frutas.txt", "w")
archivo.writelines(frutas)
archivo.close()

# Verificar contenido
archivo = open("frutas.txt", "r")
print(archivo.read())
archivo.close()
\end{lstlisting}

\vspace{-0.2cm}
\begin{tcolorbox}[colback=green!10,colframe=green,title={\small\faArrowCircleRight \ Salida},coltitle=fg0,left=3pt,right=3pt,top=2pt,bottom=2pt,boxsep=1pt]
\small\texttt{Manzana\\
Banana\\
Cereza}
\end{tcolorbox}

\vspace{-0.1cm}
\begin{notebox}
\small\faInfoCircle\ \texttt{writelines()} no añade saltos de línea automáticamente
\end{notebox}
\end{frame}

\begin{frame}[fragile]{Modo append: Añadir sin borrar \faPlus}
\begin{lstlisting}
# Añadir al final del archivo existente
archivo = open("frutas.txt", "a")
archivo.write("Durazno\n")
archivo.write("Fresa\n")
archivo.close()

# Leer resultado
archivo = open("frutas.txt", "r")
print(archivo.read())
archivo.close()
\end{lstlisting}

\vspace{-0.2cm}
\begin{tcolorbox}[colback=green!10,colframe=green,title={\small\faArrowCircleRight \ Salida},coltitle=fg0,left=3pt,right=3pt,top=2pt,bottom=2pt,boxsep=1pt]
\small\texttt{Manzana\\
Banana\\
Cereza\\
Durazno\\
Fresa}
\end{tcolorbox}

\vspace{-0.1cm}
\begin{notebox}
\small\faCheckCircle\ Modo 'a' preserva el contenido existente
\end{notebox}
\end{frame}

\section{Bloque with}

\begin{frame}[fragile]{with: Manejo automático de archivos \faShieldAlt}
\begin{conceptbox}{¿Por qué usar with?}
Cierra automáticamente el archivo, incluso si hay errores
\end{conceptbox}

\begin{lstlisting}
# Forma tradicional (manual)
archivo = open("datos.txt", "r")
contenido = archivo.read()
archivo.close()  # Fácil olvidar

# Con with (recomendado)
with open("datos.txt", "r") as archivo:
    contenido = archivo.read()
# Archivo cerrado automáticamente aquí

print(contenido)
\end{lstlisting}

\vspace{-0.1cm}
\begin{notebox}
\small\faCheckCircle\ \texttt{with} garantiza que el archivo se cierre correctamente
\end{notebox}
\end{frame}

\begin{frame}[fragile]{with: Ejemplos prácticos \faCode}
\textbf{Lectura:}
\begin{lstlisting}
with open("entrada.txt", "r") as archivo:
    for linea in archivo:
        print(linea.strip())
\end{lstlisting}

\vspace{0.2cm}
\textbf{Escritura:}
\begin{lstlisting}
numeros = [1, 2, 3, 4, 5]

with open("numeros.txt", "w") as archivo:
    for num in numeros:
        archivo.write(f"{num}\n")
\end{lstlisting}

\vspace{-0.1cm}
\begin{notebox}
\small\faLightbulb\ Usa \texttt{with} siempre que sea posible
\end{notebox}
\end{frame}

\begin{frame}[fragile]{with: Múltiples archivos \faFolderOpen}
\begin{lstlisting}
# Abrir múltiples archivos simultáneamente
with open("entrada.txt", "r") as entrada, \
     open("salida.txt", "w") as salida:
    
    for linea in entrada:
        # Convertir a mayúsculas
        salida.write(linea.upper())

print("Archivo procesado")
\end{lstlisting}

\vspace{-0.2cm}
\begin{tcolorbox}[colback=green!10,colframe=green,title={\small\faArrowCircleRight \ Salida},coltitle=fg0,left=3pt,right=3pt,top=2pt,bottom=2pt,boxsep=1pt]
\small\texttt{Archivo procesado}
\end{tcolorbox}

\vspace{-0.1cm}
\begin{notebox}
\small\faInfoCircle\ Ambos archivos se cierran automáticamente
\end{notebox}
\end{frame}

\section{Archivos CSV}

\begin{frame}[fragile]{CSV: Comma-Separated Values \faTable}
\begin{conceptbox}{¿Qué es CSV?}
Formato de texto plano para datos tabulares, separados por comas
\end{conceptbox}

\textbf{Ejemplo de archivo CSV:} \texttt{estudiantes.csv}
\begin{tcolorbox}[colback=bg1,colframe=brown,left=3pt,right=3pt,top=2pt,bottom=2pt]
\small\texttt{nombre,edad,nota\\
Ana,20,85\\
Carlos,22,90\\
María,21,88}
\end{tcolorbox}

\vspace{0.2cm}
\textbf{Ventajas:}
\small
\begin{itemize}
    \item[\faCheck] Simple y universal
    \item[\faCheck] Compatible con Excel, Google Sheets, etc.
    \item[\faCheck] Fácil de procesar programáticamente
\end{itemize}
\end{frame}

\begin{frame}[fragile]{Leer CSV manualmente \faFileAlt}
\begin{lstlisting}
# Leer CSV sin módulo csv
with open("estudiantes.csv", "r") as archivo:
    # Leer encabezado
    encabezado = archivo.readline().strip().split(",")
    print(f"Columnas: {encabezado}")
    
    # Leer datos
    for linea in archivo:
        datos = linea.strip().split(",")
        nombre, edad, nota = datos
        print(f"{nombre}: {nota} puntos")
\end{lstlisting}

\vspace{-0.2cm}
\begin{tcolorbox}[colback=green!10,colframe=green,title={\small\faArrowCircleRight \ Salida},coltitle=fg0,left=3pt,right=3pt,top=2pt,bottom=2pt,boxsep=1pt]
\small\texttt{Columnas: ['nombre', 'edad', 'nota']\\
Ana: 85 puntos\\
Carlos: 90 puntos\\
María: 88 puntos}
\end{tcolorbox}
\end{frame}

\section{Módulo csv}

\begin{frame}[fragile]{Módulo csv: Lectura con reader \faBookOpen}
\begin{lstlisting}
import csv

# Leer CSV con el módulo csv
with open("estudiantes.csv", "r") as archivo:
    lector = csv.reader(archivo)
    
    # Obtener encabezado
    encabezado = next(lector)
    print(f"Columnas: {encabezado}")
    
    # Leer filas
    for fila in lector:
        nombre, edad, nota = fila
        print(f"{nombre} ({edad} años): {nota}")
\end{lstlisting}

\vspace{-0.2cm}
\begin{tcolorbox}[colback=green!10,colframe=green,title={\small\faArrowCircleRight \ Salida},coltitle=fg0,left=3pt,right=3pt,top=2pt,bottom=2pt,boxsep=1pt]
\small\texttt{Columnas: ['nombre', 'edad', 'nota']\\
Ana (20 años): 85\\
Carlos (22 años): 90\\
María (21 años): 88}
\end{tcolorbox}
\end{frame}

\begin{frame}[fragile]{Módulo csv: DictReader \faBook}
\begin{lstlisting}
import csv

# Leer CSV como diccionarios
with open("estudiantes.csv", "r") as archivo:
    lector = csv.DictReader(archivo)
    
    for fila in lector:
        print(f"Nombre: {fila['nombre']}")
        print(f"Edad: {fila['edad']}")
        print(f"Nota: {fila['nota']}")
        print("---")
\end{lstlisting}

\vspace{-0.2cm}
\begin{tcolorbox}[colback=green!10,colframe=green,title={\small\faArrowCircleRight \ Salida},coltitle=fg0,left=3pt,right=3pt,top=2pt,bottom=2pt,boxsep=1pt]
\small\texttt{Nombre: Ana\\
Edad: 20\\
Nota: 85\\
---\\
...}
\end{tcolorbox}

\vspace{-0.1cm}
\begin{notebox}
\small\faLightbulb\ \texttt{DictReader} es más legible y mantenible
\end{notebox}
\end{frame}

\begin{frame}[fragile]{Módulo csv: Escritura con writer \faPencilAlt}
\begin{lstlisting}
import csv

# Escribir datos en CSV
estudiantes = [
    ["nombre", "edad", "nota"],
    ["Pedro", 23, 92],
    ["Laura", 20, 87],
    ["Juan", 21, 89]
]

with open("nuevo_estudiantes.csv", "w", newline="") as archivo:
    escritor = csv.writer(archivo)
    escritor.writerows(estudiantes)

print("CSV creado exitosamente")
\end{lstlisting}

\vspace{-0.2cm}
\begin{tcolorbox}[colback=green!10,colframe=green,title={\small\faArrowCircleRight \ Salida},coltitle=fg0,left=3pt,right=3pt,top=2pt,bottom=2pt,boxsep=1pt]
\small\texttt{CSV creado exitosamente}
\end{tcolorbox}

\vspace{-0.1cm}
\begin{notebox}
\small\faInfoCircle\ \texttt{newline=""} evita líneas en blanco en Windows
\end{notebox}
\end{frame}

\begin{frame}[fragile]{Módulo csv: DictWriter \faEdit}
\begin{lstlisting}
import csv

# Escribir diccionarios como CSV
estudiantes = [
    {"nombre": "Sofia", "edad": 22, "nota": 95},
    {"nombre": "Diego", "edad": 20, "nota": 88},
]

with open("estudiantes_dict.csv", "w", newline="") as archivo:
    campos = ["nombre", "edad", "nota"]
    escritor = csv.DictWriter(archivo, fieldnames=campos)
    
    escritor.writeheader()  # Escribir encabezado
    escritor.writerows(estudiantes)

print("CSV con diccionarios creado")
\end{lstlisting}

\vspace{-0.2cm}
\begin{tcolorbox}[colback=green!10,colframe=green,title={\small\faArrowCircleRight \ Salida},coltitle=fg0,left=3pt,right=3pt,top=2pt,bottom=2pt,boxsep=1pt]
\small\texttt{CSV con diccionarios creado}
\end{tcolorbox}
\end{frame}

\section{Rutas de archivos}

\begin{frame}[fragile]{Trabajo con rutas: Módulo os.path \faRoute}
\begin{lstlisting}
import os

# Ruta actual
directorio = os.getcwd()
print(f"Directorio actual: {directorio}")

# Verificar existencia
existe = os.path.exists("datos.txt")
print(f"¿Existe datos.txt?: {existe}")

# Verificar si es archivo o directorio
if os.path.exists("datos.txt"):
    es_archivo = os.path.isfile("datos.txt")
    es_dir = os.path.isdir("datos.txt")
    print(f"¿Es archivo?: {es_archivo}")
    print(f"¿Es directorio?: {es_dir}")
\end{lstlisting}

\vspace{-0.1cm}
\begin{notebox}
\small\faCheckCircle\ Siempre verifica la existencia antes de abrir
\end{notebox}
\end{frame}

\begin{frame}[fragile]{Construir rutas multiplataforma \faFolderTree}
\begin{lstlisting}
import os

# Construir ruta correctamente
carpeta = "datos"
archivo = "estudiantes.csv"

# Forma incorrecta (solo funciona en algunos SO)
# ruta = "datos/estudiantes.csv"

# Forma correcta (funciona en todos los SO)
ruta = os.path.join(carpeta, archivo)
print(f"Ruta: {ruta}")

# Obtener información
nombre = os.path.basename(ruta)
directorio = os.path.dirname(ruta)

print(f"Nombre: {nombre}")
print(f"Directorio: {directorio}")
\end{lstlisting}

\vspace{-0.1cm}
\begin{notebox}
\small\faLightbulb\ \texttt{os.path.join()} usa el separador correcto del SO
\end{notebox}
\end{frame}

\begin{frame}[fragile]{Módulo pathlib (moderno) \faFileCode}
\begin{lstlisting}
from pathlib import Path

# Crear objeto Path
ruta = Path("datos") / "estudiantes.csv"
print(f"Ruta: {ruta}")

# Verificar existencia
if ruta.exists():
    print("El archivo existe")
    
    # Leer contenido
    contenido = ruta.read_text()
    print(contenido[:50])  # Primeros 50 caracteres

# Obtener partes
print(f"Nombre: {ruta.name}")
print(f"Sufijo: {ruta.suffix}")
print(f"Padre: {ruta.parent}")
\end{lstlisting}

\vspace{-0.1cm}
\begin{notebox}
\small\faInfoCircle\ \texttt{pathlib} es más moderno y orientado a objetos
\end{notebox}
\end{frame}

\section{Manejo de errores}

\begin{frame}[fragile]{Manejo de errores con archivos \faExclamationTriangle}
\begin{lstlisting}
# Manejo robusto de archivos
try:
    with open("archivo_inexistente.txt", "r") as archivo:
        contenido = archivo.read()
        print(contenido)
except FileNotFoundError:
    print("Error: El archivo no existe")
except PermissionError:
    print("Error: No tienes permiso para leer el archivo")
except Exception as e:
    print(f"Error inesperado: {e}")
\end{lstlisting}

\vspace{-0.2cm}
\begin{tcolorbox}[colback=red!10,colframe=red,title={\small\faExclamationTriangle \ Salida},coltitle=fg0,left=3pt,right=3pt,top=2pt,bottom=2pt,boxsep=1pt]
\small\texttt{Error: El archivo no existe}
\end{tcolorbox}

\vspace{-0.1cm}
\begin{notebox}
\small\faShieldAlt\ Siempre maneja posibles errores con archivos
\end{notebox}
\end{frame}

\section{Ejemplo integrador}

\begin{frame}[fragile]{Proyecto: Agenda de contactos (1/5) \faAddressBook}
\textbf{Estructura del programa:}
\small
\begin{itemize}
    \item[\faFile] Archivo CSV con contactos
    \item[\faPlus] Agregar nuevos contactos
    \item[\faList] Listar todos los contactos
    \item[\faSearch] Buscar contacto por nombre
    \item[\faTrash] Eliminar contacto
\end{itemize}

\vspace{0.2cm}
\textbf{Archivo:} \texttt{agenda.py}
\begin{lstlisting}
import csv
import os

ARCHIVO_CONTACTOS = "contactos.csv"
\end{lstlisting}
\end{frame}

\begin{frame}[fragile]{Proyecto: Agenda de contactos (2/5) \faCode}
\begin{lstlisting}
def inicializar_archivo():
    """Crea el archivo si no existe"""
    if not os.path.exists(ARCHIVO_CONTACTOS):
        with open(ARCHIVO_CONTACTOS, "w", newline="") as archivo:
            escritor = csv.writer(archivo)
            escritor.writerow(["nombre", "telefono", "email"])

def agregar_contacto(nombre, telefono, email):
    """Agrega un nuevo contacto"""
    with open(ARCHIVO_CONTACTOS, "a", newline="") as archivo:
        escritor = csv.writer(archivo)
        escritor.writerow([nombre, telefono, email])
    print(f"Contacto {nombre} agregado exitosamente")
\end{lstlisting}
\end{frame}

\begin{frame}[fragile]{Proyecto: Agenda de contactos (3/5) \faList}
\begin{lstlisting}
def listar_contactos():
    """Muestra todos los contactos"""
    with open(ARCHIVO_CONTACTOS, "r") as archivo:
        lector = csv.DictReader(archivo)
        print("\n--- CONTACTOS ---")
        for i, contacto in enumerate(lector, 1):
            print(f"{i}. {contacto['nombre']}")
            print(f"   Tel: {contacto['telefono']}")
            print(f"   Email: {contacto['email']}")
            print()
\end{lstlisting}
\end{frame}

\begin{frame}[fragile]{Proyecto: Agenda de contactos (4/5) \faSearch}
\begin{lstlisting}
def buscar_contacto(nombre_buscar):
    """Busca un contacto por nombre"""
    with open(ARCHIVO_CONTACTOS, "r") as archivo:
        lector = csv.DictReader(archivo)
        encontrado = False
        
        for contacto in lector:
            if nombre_buscar.lower() in contacto['nombre'].lower():
                print(f"\nNombre: {contacto['nombre']}")
                print(f"Teléfono: {contacto['telefono']}")
                print(f"Email: {contacto['email']}")
                encontrado = True
        
        if not encontrado:
            print(f"No se encontró: {nombre_buscar}")
\end{lstlisting}
\end{frame}

\begin{frame}[fragile]{Proyecto: Agenda de contactos (5/5) \faPlayCircle}
\begin{lstlisting}
def menu():
    """Menú principal"""
    inicializar_archivo()
    
    while True:
        print("\n=== AGENDA DE CONTACTOS ===")
        print("1. Agregar contacto")
        print("2. Listar contactos")
        print("3. Buscar contacto")
        print("4. Salir")
        
        opcion = input("\nElige opción: ")
        
        if opcion == "1":
            nombre = input("Nombre: ")
            telefono = input("Teléfono: ")
            email = input("Email: ")
            agregar_contacto(nombre, telefono, email)
\end{lstlisting}

\small(Continuación en siguiente slide...)
\end{frame}

\begin{frame}[fragile]{Proyecto: Agenda - Ejecución \faTerminal}
\begin{lstlisting}
        elif opcion == "2":
            listar_contactos()
        elif opcion == "3":
            nombre = input("Buscar por nombre: ")
            buscar_contacto(nombre)
        elif opcion == "4":
            print("¡Hasta luego!")
            break
        else:
            print("Opción inválida")

# Ejecutar programa
if __name__ == "__main__":
    menu()
\end{lstlisting}

\vspace{-0.1cm}
\begin{notebox}
\small\faRocket\ ¡Programa completo y funcional con CSV!
\end{notebox}
\end{frame}

\begin{frame}[fragile]{Proyecto alternativo: Analizador CSV \faChartBar}
\begin{lstlisting}
import csv

def analizar_ventas(archivo_csv):
    """Analiza datos de ventas desde CSV"""
    total = 0
    productos = []
    
    with open(archivo_csv, "r") as archivo:
        lector = csv.DictReader(archivo)
        
        for fila in lector:
            producto = fila['producto']
            precio = float(fila['precio'])
            cantidad = int(fila['cantidad'])
            
            subtotal = precio * cantidad
            total += subtotal
            productos.append((producto, subtotal))
\end{lstlisting}

\small(Continuación en siguiente slide...)
\end{frame}

\begin{frame}[fragile]{Proyecto alternativo: Analizador (cont.) \faFileAlt}
\begin{lstlisting}
    # Mostrar resultados
    print("\n=== ANÁLISIS DE VENTAS ===")
    print(f"Total vendido: ${total:.2f}\n")
    
    # Top productos
    productos.sort(key=lambda x: x[1], reverse=True)
    print("Top 3 productos:")
    for i, (producto, monto) in enumerate(productos[:3], 1):
        print(f"{i}. {producto}: ${monto:.2f}")

# Usar
analizar_ventas("ventas.csv")
\end{lstlisting}

\vspace{-0.1cm}
\begin{notebox}
\small\faChartLine\ Procesar y analizar datos es muy común en Python
\end{notebox}
\end{frame}

\begin{frame}{Buenas prácticas con archivos \faCheckDouble}
\begin{conceptbox}{Recomendaciones}
\small
\begin{itemize}
    \item[\faCheck] Usa siempre \texttt{with} para abrir archivos
    \item[\faCheck] Verifica la existencia con \texttt{os.path.exists()}
    \item[\faCheck] Maneja errores con \texttt{try-except}
    \item[\faCheck] Usa \texttt{os.path.join()} para rutas multiplataforma
    \item[\faCheck] Cierra archivos explícitamente si no usas \texttt{with}
    \item[\faCheck] Para CSV, usa el módulo \texttt{csv}
    \item[\faCheck] Valida datos antes de escribir
    \item[\faCheck] Usa codificación UTF-8 si trabajas con acentos
\end{itemize}
\end{conceptbox}
\end{frame}

\begin{frame}[fragile]{Codificación de caracteres \faLanguage}
\begin{conceptbox}{UTF-8 para caracteres especiales}
Importante cuando trabajas con acentos, ñ, etc.
\end{conceptbox}

\begin{lstlisting}
# Especificar codificación UTF-8
with open("texto.txt", "w", encoding="utf-8") as archivo:
    archivo.write("Año: 2024\n")
    archivo.write("Niño español\n")
    archivo.write("Caña, piñata\n")

# Leer con UTF-8
with open("texto.txt", "r", encoding="utf-8") as archivo:
    print(archivo.read())
\end{lstlisting}

\vspace{-0.1cm}
\begin{notebox}
\small\faExclamationTriangle\ Usa \texttt{encoding="utf-8"} para evitar problemas
\end{notebox}
\end{frame}

\begin{frame}{Resumen \faListUl}
\begin{conceptbox}{¿Qué aprendimos hoy?}
\small
\begin{itemize}
    \item[\faCheck] Función \texttt{open()} y sus modos
    \item[\faCheck] Lectura: \texttt{read()}, \texttt{readline()}, \texttt{readlines()}
    \item[\faCheck] Escritura: \texttt{write()}, \texttt{writelines()}
    \item[\faCheck] Bloque \texttt{with} para manejo automático
    \item[\faCheck] Archivos CSV con módulo \texttt{csv}
    \item[\faCheck] Rutas con \texttt{os.path} y \texttt{pathlib}
    \item[\faCheck] Manejo de errores
    \item[\faCheck] Proyecto: Agenda de contactos
\end{itemize}
\end{conceptbox}

\vspace{0.3cm}
\begin{center}
\large\textcolor{brown}{\faFileAlt\ ¡Los archivos dan persistencia a tus programas!}
\end{center}
\end{frame}

\begin{frame}
\begin{center}
\Huge\textcolor{brown}{\faQuestionCircle}\\
\vspace{0.5cm}
\Huge ¡Preguntas!\\
\vspace{1cm}
\Large\textcolor{brownbright}{Gracias por su atención}
\end{center}
\end{frame}

\end{document}
